\thispagestyle{plain}
\noindent{\sffamily\heiti\Large KSSOLV Toolbox 中文用户手册}

\vspace{5mm}
\noindent 独立应用程序、MATLAB\regmark{} 工具箱与\ksonline{}\par
\noindent 适用版本：V0.5.0\quad 编写日期：2026 年 10 月 21 日

\vspace{7mm}
本书介绍 KSSOLV Toolbox 的安装、结构准备、电子结构计算和结果分析。教程从小型分子与晶体出发，逐步延伸到表面、异质结、扭转结构、低维模型与软物质。本书适用于独立应用程序、MATLAB\regmark{} 工具箱和\ksonline{}三种使用方式，共通的建模与计算步骤统一讲解，安装、文件管理和远程运行的差别在相应章节说明。

本书界面名称按中文界面书写；仍以英文显示的字段保留原名。软件更新后，部分窗口的位置和外观可能有所调整。计算案例的输入条件和参考结果一并给出，便于读者理解结果的来源。

\noindent{\sffamily\heiti 品牌}

瀚海量子\regmark{}标识和 KSSOLV Toolbox 名称用于标明本产品。教程和插图说明围绕 KSSOLV Toolbox 编写。MATLAB\regmark{} 是 The MathWorks, Inc. 的注册商标。

\noindent{\sffamily\heiti 封面}

封面晶体取自随书硅案例，沿 $[100]$ 方向延展三个常规晶胞。球体表示硅原子，连线连接第一近邻；青绿色的四条键突出金刚石结构中的四面体配位。

\vfill
\noindent{\small © 2026 合肥瀚海量子科技有限公司\par
\noindent Hefei Hanhai Quantum Technology Co., Ltd.}
\clearpage

\noindent{\sffamily\heiti\Large 关于本书}

\vspace{4mm}
\noindent{\sffamily\heiti 版本与编写依据}

本书面向即将发布的 KSSOLV Toolbox V0.5.0，界面和功能依据编写期间的工作构建。安装包、运行环境和在线版部署可分别更新，使用时请在\ui{主页 → 帮助 → 关于 KSSOLV}中确认版本。

\noindent{\sffamily\heiti 建议的阅读顺序}

第一次使用时，可先在第\ref{ch:introduction}章选择使用方式，再按安装或在线版章节准备运行环境，随后阅读第\ref{ch:workspace}章熟悉工作区。第\ref{ch:first}章的氢分子教程带读者完成一次从结构到结果的计算，可按步骤操作并核对每一阶段的结果。

准备自己的材料模型时，可按第\ref{ch:modeling}章编辑原子结构，再到第\ref{ch:advanced}章学习材料实例；第\ref{ch:modeling-tools}章提供各项建模工具的查阅说明。计算参数一章解释常见设置的含义；工作流、体数据与结果章节则用于进一步安排任务和阅读输出。不必在第一次操作前读完所有章节。

随书的 \file{examples/} 目录提供结构、项目、复算脚本与体数据。练习时从指定文件开始，保存自己的项目副本；比较数值时，同时核对结构、赝势和计算条件。

\noindent{\sffamily\heiti 获取软件与帮助}

\begin{itemize}
\item 软件与发行包：\url{https://github.com/USTC-Hu-Group/KSSOLV-Toolbox}
\item 在线版：\ksonline{} 的瀚海量子\regmark{}官方在线版即将推出，并支持集群私有化部署。
\item 问题反馈：上述仓库的 Issues 页面；也可从应用\ui{帮助}菜单进入相关资源。
\end{itemize}

\clearpage
